\documentclass[12pt,a4paper]{article}

\usepackage[T2A]{fontenc}
\usepackage[utf8]{inputenc}
\usepackage[russian]{babel}
\usepackage{amsmath,amssymb}
\usepackage{enumitem}
\usepackage{array}
\usepackage{longtable}
\usepackage{listings}
\usepackage{xcolor}
\usepackage{geometry}
\usepackage{tikz}
\usetikzlibrary{arrows.meta,positioning}

\geometry{left=2.2cm,right=2.2cm,top=2cm,bottom=2cm}

\lstset{
    basicstyle=\ttfamily\small,
    keywordstyle=\color{blue},
    commentstyle=\color{gray},
    stringstyle=\color{teal},
    showstringspaces=false,
    frame=single,
    breaklines=true
}

\title{Билет №66 \\ \small Параллельные процессы, схемы порождения и управления. Организация межпроцессного взаимодействия: общая память, обмен сообщениями, организация почтовых ящиков. Модели согласованности данных. (ИТМО)}
\author{Сериков Владислав Александрович}
\date{16 мая 2026}

\begin{document}

\maketitle
\tableofcontents
\newpage

\section{Параллельные процессы, порождение, управление и межпроцессное взаимодействие}

\subsection{Базовые понятия}

\textbf{Процесс} --- это выполняющаяся программа вместе со своим адресным пространством, состоянием процессора, открытыми файлами, правами доступа и другими ресурсами операционной системы.

\textbf{Поток выполнения} --- минимальная единица планирования, которая исполняется внутри процесса. Потоки одного процесса обычно разделяют адресное пространство и ресурсы процесса, но имеют собственные регистры, стек и состояние выполнения.

\textbf{Параллельные процессы} --- процессы или потоки, выполнение которых перекрывается во времени. В зависимости от аппаратной платформы различают:
\begin{itemize}
    \item \textbf{конкурентное выполнение} --- процессы логически выполняются одновременно, но физически могут разделять один процессор;
    \item \textbf{параллельное выполнение} --- процессы действительно исполняются одновременно на разных процессорных ядрах;
    \item \textbf{распределённое выполнение} --- процессы выполняются на разных узлах сети и взаимодействуют через коммуникационную среду.
\end{itemize}

Параллельные процессы могут быть:
\begin{itemize}
    \item \textbf{независимыми}, если их выполнение не влияет друг на друга;
    \item \textbf{взаимодействующими}, если они обмениваются данными или синхронизируют свои действия;
    \item \textbf{конкурирующими}, если они претендуют на одни и те же ресурсы.
\end{itemize}

\subsection{Причины использования параллельных процессов}

Параллельные процессы используются для:
\begin{itemize}
    \item повышения производительности за счёт использования нескольких ядер;
    \item повышения отзывчивости интерактивных систем;
    \item естественного моделирования независимых сущностей, например клиентов сервера;
    \item разделения вычислений на независимые подзадачи;
    \item организации асинхронной обработки ввода-вывода;
    \item построения распределённых систем.
\end{itemize}

\subsection{Основные проблемы параллельного выполнения}

Параллельность порождает ряд фундаментальных проблем:
\begin{itemize}
    \item \textbf{состояния гонки} --- результат зависит от порядка выполнения операций;
    \item \textbf{критические секции} --- участки кода, работающие с разделяемыми ресурсами;
    \item \textbf{взаимная блокировка} --- процессы бесконечно ждут ресурсы друг друга;
    \item \textbf{голодание} --- процесс бесконечно откладывается планировщиком или синхронизационным механизмом;
    \item \textbf{инверсия приоритетов} --- низкоприоритетный процесс удерживает ресурс, нужный высокоприоритетному;
    \item \textbf{недетерминизм} --- разные запуски программы могут давать разные допустимые результаты.
\end{itemize}

\subsection{Состояние гонки и критическая секция}

\textbf{Состояние гонки} возникает, когда несколько процессов одновременно обращаются к разделяемым данным, по крайней мере один из них выполняет запись, а итоговый результат зависит от относительного порядка операций.

Например, пусть два процесса увеличивают общую переменную:
\[
x := x + 1.
\]
На машинном уровне операция может состоять из трёх шагов:
\[
\text{read } x,\quad \text{increment},\quad \text{write } x.
\]

Если изначально $x = 0$, то возможна следующая история:

\[
\begin{array}{c|c|c}
\text{Шаг} & P_1 & P_2 \\
\hline
1 & r_1 := x = 0 & \\
2 & & r_2 := x = 0 \\
3 & r_1 := r_1 + 1 = 1 & \\
4 & & r_2 := r_2 + 1 = 1 \\
5 & x := r_1 = 1 & \\
6 & & x := r_2 = 1
\end{array}
\]

Ожидаемый результат после двух инкрементов равен $2$, но фактически получено $1$.

\textbf{Критическая секция} --- участок программы, в котором процесс обращается к разделяемому ресурсу, требующему согласованного доступа.

Корректное решение задачи критической секции должно удовлетворять следующим свойствам:
\begin{enumerate}
    \item \textbf{Взаимное исключение}: одновременно в критической секции находится не более одного процесса.
    \item \textbf{Прогресс}: если критическая секция свободна, то выбор следующего процесса не должен откладываться бесконечно.
    \item \textbf{Ограниченное ожидание}: существует верхняя граница числа входов других процессов в критическую секцию после того, как данный процесс запросил вход.
\end{enumerate}

\subsection{Схемы порождения параллельных процессов}

Порождение процессов может организовываться несколькими схемами.

\subsubsection{Статическое порождение}

При \textbf{статическом порождении} все процессы создаются заранее, до начала основной фазы вычислений. Число процессов фиксировано.

Типичная схема:
\[
P_0 \rightarrow \{P_1, P_2, \ldots, P_n\}.
\]

Пример:
\begin{itemize}
    \item программа заранее создаёт пул из $n$ рабочих потоков;
    \item каждый поток ожидает задачи из очереди;
    \item после завершения всех задач потоки уничтожаются.
\end{itemize}

Преимущества:
\begin{itemize}
    \item простота управления;
    \item отсутствие постоянных накладных расходов на создание процессов;
    \item предсказуемое потребление ресурсов.
\end{itemize}

Недостатки:
\begin{itemize}
    \item плохо подходит для динамически изменяющейся нагрузки;
    \item может приводить к недоиспользованию ресурсов или перегрузке.
\end{itemize}

\subsubsection{Динамическое порождение}

При \textbf{динамическом порождении} процессы создаются по мере необходимости.

Например, сервер может создавать отдельный процесс или поток для каждого клиента:
\[
\text{client request} \Rightarrow \text{create worker}.
\]

Преимущества:
\begin{itemize}
    \item гибкость;
    \item адаптация к текущей нагрузке;
    \item естественная модель для серверных приложений.
\end{itemize}

Недостатки:
\begin{itemize}
    \item большие накладные расходы при частом создании;
    \item риск исчерпания ресурсов;
    \item сложность контроля числа процессов.
\end{itemize}

\subsubsection{Иерархическое порождение}

Процессы могут образовывать дерево:

\[
P_0
\rightarrow
\{P_1, P_2\},
\quad
P_1 \rightarrow \{P_3, P_4\},
\quad
P_2 \rightarrow \{P_5\}.
\]

Такую модель используют многие операционные системы. Например, в UNIX-подобных системах процесс создаётся системным вызовом \texttt{fork()}, после чего может выполнить другую программу через \texttt{exec()}.

\subsubsection{Модель fork--join}

Одна из классических схем параллельного программирования --- \textbf{fork--join}.

\[
\text{fork}: \text{ разделение вычисления на параллельные ветви},
\]
\[
\text{join}: \text{ ожидание завершения ветвей и объединение результатов}.
\]

Общий вид:

\[
\begin{array}{c}
\text{основной процесс} \\
\downarrow \\
\text{fork} \\
\swarrow \quad \searrow \\
P_1 \quad P_2 \\
\searrow \quad \swarrow \\
\text{join} \\
\downarrow \\
\text{продолжение}
\end{array}
\]

Минимальный пример: параллельное суммирование массива.

Пусть имеется массив:
\[
A = [a_1, a_2, \ldots, a_n].
\]

Алгоритм:
\begin{enumerate}
    \item Разделить массив на две части:
    \[
    A_L = [a_1,\ldots,a_{n/2}], \quad A_R = [a_{n/2+1},\ldots,a_n].
    \]
    \item Создать два процесса:
    \[
    P_L \text{ считает сумму } A_L,\quad P_R \text{ считает сумму } A_R.
    \]
    \item Дождаться завершения обоих процессов.
    \item Сложить частичные суммы:
    \[
    S = S_L + S_R.
    \]
\end{enumerate}

В псевдокоде:
\[
\begin{array}{l}
\textbf{function } parallel\_sum(A):\\
\quad \textbf{if } |A| \leq T \textbf{ then return } sequential\_sum(A)\\
\quad A_L, A_R := split(A)\\
\quad f_L := fork(parallel\_sum(A_L))\\
\quad f_R := fork(parallel\_sum(A_R))\\
\quad S_L := join(f_L)\\
\quad S_R := join(f_R)\\
\quad \textbf{return } S_L + S_R
\end{array}
\]

Здесь $T$ --- порог, после которого выгоднее считать последовательно.

\subsubsection{Модель master--worker}

В модели \textbf{master--worker} один управляющий процесс распределяет задачи между рабочими процессами.

\[
\text{Master} \rightarrow \{W_1, W_2, \ldots, W_k\}.
\]

Шаги:
\begin{enumerate}
    \item Master формирует очередь задач.
    \item Worker запрашивает задачу.
    \item Master выдаёт задачу свободному worker'у.
    \item Worker выполняет задачу и возвращает результат.
    \item Master агрегирует результаты.
\end{enumerate}

Минимальный пример: обработка набора файлов.
\[
\text{файлы } F_1,\ldots,F_m
\]
распределяются между рабочими процессами:
\[
W_i \leftarrow F_j.
\]

Преимущества:
\begin{itemize}
    \item простая балансировка нагрузки;
    \item удобство централизованного контроля;
    \item масштабируемость до разумного числа работников.
\end{itemize}

Недостатки:
\begin{itemize}
    \item master может стать узким местом;
    \item отказ master'а может нарушить работу всей системы.
\end{itemize}

\subsubsection{Конвейерная схема}

В \textbf{конвейере} обработка разбивается на последовательность стадий:

\[
P_1 \rightarrow P_2 \rightarrow \cdots \rightarrow P_k.
\]

Каждая стадия выполняется отдельным процессом. Данные проходят через стадии по очереди.

Пример компилятора:
\[
\text{лексический анализ}
\rightarrow
\text{синтаксический анализ}
\rightarrow
\text{семантический анализ}
\rightarrow
\text{генерация кода}.
\]

Преимущество конвейера состоит в том, что разные элементы данных одновременно находятся на разных стадиях обработки.

\subsection{Схемы управления параллельными процессами}

\subsubsection{Централизованное управление}

При централизованном управлении один процесс принимает основные решения:
\begin{itemize}
    \item кому выдать задачу;
    \item когда остановить вычисление;
    \item как агрегировать результаты;
    \item как обрабатывать ошибки.
\end{itemize}

Примером является модель master--worker.

Плюсы:
\begin{itemize}
    \item простота проектирования;
    \item глобальный контроль состояния;
    \item удобная диагностика.
\end{itemize}

Минусы:
\begin{itemize}
    \item единая точка отказа;
    \item ограниченная масштабируемость;
    \item возможное узкое место.
\end{itemize}

\subsubsection{Децентрализованное управление}

При децентрализованном управлении процессы сами принимают решения на основе локальной информации и сообщений.

Примеры:
\begin{itemize}
    \item одноранговые сети;
    \item распределённые хеш-таблицы;
    \item алгоритмы консенсуса;
    \item распределённая балансировка нагрузки.
\end{itemize}

Плюсы:
\begin{itemize}
    \item отсутствие единой точки отказа;
    \item высокая масштабируемость;
    \item естественная пригодность для распределённых систем.
\end{itemize}

Минусы:
\begin{itemize}
    \item сложность доказательства корректности;
    \item накладные расходы на координацию;
    \item необходимость учитывать задержки, отказы и рассинхронизацию часов.
\end{itemize}

\subsubsection{Синхронное и асинхронное управление}

При \textbf{синхронном} управлении процессы выполняют шаги согласованно, например по тактам или барьерам.

При \textbf{асинхронном} управлении каждый процесс движется независимо, а взаимодействие происходит через события, сообщения или разделяемые структуры.

Асинхронные системы сложнее анализировать, поскольку порядок событий заранее неизвестен.

\subsection{Организация межпроцессного взаимодействия}

\textbf{Межпроцессное взаимодействие} --- это совокупность механизмов, позволяющих процессам обмениваться данными и синхронизировать выполнение.

Основные модели:
\begin{enumerate}
    \item общая память;
    \item обмен сообщениями;
    \item почтовые ящики;
    \item каналы;
    \item сигналы;
    \item сокеты;
    \item удалённый вызов процедур;
    \item разделяемые файлы и базы данных.
\end{enumerate}

Наиболее фундаментальными считаются две модели:
\begin{itemize}
    \item \textbf{shared memory} --- общая память;
    \item \textbf{message passing} --- обмен сообщениями.
\end{itemize}

\section{Межпроцессное взаимодействие через общую память}

\subsection{Общая память}

\textbf{Общая память} --- область памяти, доступная нескольким процессам или потокам.

В многопоточной программе общая память обычно возникает естественно: все потоки одного процесса имеют доступ к глобальным переменным и куче.

В многопроцессной системе общая память может создаваться явно средствами ОС, например:
\begin{itemize}
    \item POSIX shared memory;
    \item System V shared memory;
    \item memory-mapped files;
    \item разделяемые сегменты памяти.
\end{itemize}

\subsection{Структура взаимодействия через общую память}

Общая схема:

\[
P_1 \leftrightarrow M \leftrightarrow P_2,
\]

где $M$ --- разделяемая область памяти.

Процессы читают и записывают данные в $M$:
\[
P_i: \quad read(M), \quad write(M).
\]

\subsection{Преимущества общей памяти}

\begin{itemize}
    \item высокая скорость обмена данными;
    \item отсутствие необходимости копировать большие объёмы информации;
    \item удобство для многопоточных приложений;
    \item естественная модель для совместного доступа к структурам данных.
\end{itemize}

\subsection{Недостатки общей памяти}

\begin{itemize}
    \item необходимость явной синхронизации;
    \item возможность состояний гонки;
    \item сложность отладки;
    \item проблемы кэш-когерентности на многопроцессорных системах;
    \item сложность масштабирования на распределённые системы.
\end{itemize}

\subsection{Классическая задача производителя и потребителя}

Одной из базовых задач взаимодействия через общую память является задача \textbf{производителя и потребителя}.

Имеются:
\begin{itemize}
    \item производитель, создающий элементы данных;
    \item потребитель, обрабатывающий элементы;
    \item общий буфер размера $N$.
\end{itemize}

Производитель не должен помещать данные в полный буфер, а потребитель не должен читать из пустого буфера.

\subsubsection{Решение с семафорами}

Используются:
\begin{itemize}
    \item \texttt{mutex} --- бинарный семафор для защиты буфера;
    \item \texttt{empty} --- счётчик свободных ячеек;
    \item \texttt{full} --- счётчик занятых ячеек.
\end{itemize}

Начальные значения:
\[
mutex = 1,\quad empty = N,\quad full = 0.
\]

Производитель:
\[
\begin{array}{l}
\textbf{while true do}\\
\quad item := produce()\\
\quad wait(empty)\\
\quad wait(mutex)\\
\quad insert(item)\\
\quad signal(mutex)\\
\quad signal(full)
\end{array}
\]

Потребитель:
\[
\begin{array}{l}
\textbf{while true do}\\
\quad wait(full)\\
\quad wait(mutex)\\
\quad item := remove()\\
\quad signal(mutex)\\
\quad signal(empty)\\
\quad consume(item)
\end{array}
\]

Здесь:
\[
wait(S)
\]
уменьшает семафор $S$, если $S > 0$, иначе блокирует процесс.

\[
signal(S)
\]
увеличивает семафор $S$ и, возможно, пробуждает ожидающий процесс.

\subsubsection{Инвариант буфера}

Для корректности важно, что во все моменты времени выполняется:
\[
empty + full = N.
\]

Кроме того:
\[
0 \leq empty \leq N,
\quad
0 \leq full \leq N.
\]

\subsection{Алгоритм Петерсона для двух процессов}

Алгоритм Петерсона --- классический программный алгоритм взаимного исключения для двух процессов.

Пусть есть два процесса $P_0$ и $P_1$.

Общие переменные:
\[
flag[0], flag[1] \in \{false, true\},
\]
\[
turn \in \{0,1\}.
\]

Переменная $flag[i]$ означает, что процесс $P_i$ хочет войти в критическую секцию.

Псевдокод процесса $P_i$:
\[
\begin{array}{l}
flag[i] := true;\\
turn := 1 - i;\\
\textbf{while } flag[1-i] = true \textbf{ and } turn = 1-i \textbf{ do}\\
\quad \text{ожидать};\\
\text{critical section};\\
flag[i] := false;\\
\text{remainder section}.
\end{array}
\]

\subsubsection{Идея алгоритма}

Если оба процесса хотят войти в критическую секцию, каждый устанавливает свой флаг. Переменная \texttt{turn} отдаёт преимущество другому процессу. Тот процесс, которому не принадлежит очередь, ждёт.

\subsubsection{Теорема о взаимном исключении в алгоритме Петерсона}

\textbf{Теорема.}
В алгоритме Петерсона для двух процессов невозможно, чтобы оба процесса одновременно находились в критической секции.

\textbf{Доказательство.}

Предположим противное: процессы $P_0$ и $P_1$ одновременно находятся в критической секции.

Для входа в критическую секцию процесс $P_0$ должен был выйти из цикла:
\[
flag[1] = true \land turn = 1.
\]
Следовательно, в момент выхода $P_0$ из цикла было выполнено:
\[
flag[1] = false \quad \text{или} \quad turn \neq 1.
\]
Так как $turn$ принимает только значения $0$ и $1$, условие $turn \neq 1$ эквивалентно:
\[
turn = 0.
\]

Аналогично, для входа $P_1$ в критическую секцию он должен был выйти из цикла:
\[
flag[0] = true \land turn = 0.
\]
Следовательно, в момент выхода $P_1$ из цикла было выполнено:
\[
flag[0] = false \quad \text{или} \quad turn = 1.
\]

Поскольку оба процесса находятся в критической секции, каждый из них перед входом установил свой флаг:
\[
flag[0] = true,\quad flag[1] = true.
\]
Эти флаги сбрасываются только после выхода из критической секции. Значит, пока оба процесса находятся в критической секции, выполняется:
\[
flag[0] = true,\quad flag[1] = true.
\]

Тогда для $P_0$ единственной причиной выйти из цикла могло быть:
\[
turn = 0.
\]
Для $P_1$ единственной причиной выйти из цикла могло быть:
\[
turn = 1.
\]

Но переменная $turn$ не может одновременно иметь значения $0$ и $1$:
\[
turn = 0 \land turn = 1
\]
невозможно.

Получено противоречие. Следовательно, два процесса не могут одновременно находиться в критической секции.

Теорема доказана.
\[
\blacksquare
\]

\subsubsection{Теорема о прогрессе в алгоритме Петерсона}

\textbf{Теорема.}
Если ни один процесс не находится в критической секции и хотя бы один процесс хочет войти в неё, то выбор процесса, который войдёт в критическую секцию, не откладывается бесконечно.

\textbf{Доказательство.}

Рассмотрим возможные случаи.

\textbf{Случай 1.} Только процесс $P_i$ хочет войти в критическую секцию.

Тогда:
\[
flag[i] = true,
\quad
flag[1-i] = false.
\]
Условие ожидания процесса $P_i$:
\[
flag[1-i] = true \land turn = 1-i.
\]
Поскольку $flag[1-i] = false$, всё условие ложно. Следовательно, $P_i$ входит в критическую секцию без ожидания.

\textbf{Случай 2.} Оба процесса хотят войти в критическую секцию.

Тогда:
\[
flag[0] = flag[1] = true.
\]
Оба процесса записывают значение в переменную $turn$. Последняя запись определяет её окончательное значение.

Если:
\[
turn = 0,
\]
то процесс $P_1$ ждёт, поскольку его условие ожидания:
\[
flag[0] = true \land turn = 0
\]
истинно. Процесс $P_0$ не ждёт, поскольку его условие ожидания:
\[
flag[1] = true \land turn = 1
\]
ложно.

Если:
\[
turn = 1,
\]
то симметрично входит $P_1$, а $P_0$ ждёт.

Таким образом, в любом случае хотя бы один процесс может войти в критическую секцию. Следовательно, прогресс обеспечивается.

Теорема доказана.
\[
\blacksquare
\]

\subsection{Семафоры}

\textbf{Семафор} --- целочисленная переменная, доступ к которой осуществляется только через две атомарные операции:
\[
wait(S), \quad signal(S).
\]

Классическое определение:
\[
wait(S):
\begin{cases}
S := S - 1, & \text{если } S > 0,\\
\text{блокировка процесса}, & \text{если } S = 0.
\end{cases}
\]

\[
signal(S):
\begin{cases}
S := S + 1,\\
\text{пробуждение одного ожидающего процесса, если он есть}.
\end{cases}
\]

Различают:
\begin{itemize}
    \item \textbf{бинарный семафор}, принимающий значения $0$ и $1$;
    \item \textbf{счётный семафор}, принимающий неотрицательные целые значения.
\end{itemize}

\subsection{Мьютексы}

\textbf{Мьютекс} --- объект синхронизации для организации взаимного исключения.

Типичный шаблон:
\[
\begin{array}{l}
lock(m);\\
\text{critical section};\\
unlock(m).
\end{array}
\]

Отличия мьютекса от семафора:
\begin{itemize}
    \item мьютекс обычно имеет владельца;
    \item разблокировать мьютекс должен тот поток, который его захватил;
    \item семафор может использоваться не только для взаимного исключения, но и для подсчёта ресурсов.
\end{itemize}

\subsection{Мониторы}

\textbf{Монитор} --- высокоуровневый механизм синхронизации, объединяющий:
\begin{itemize}
    \item разделяемые данные;
    \item процедуры доступа к ним;
    \item неявное взаимное исключение;
    \item условные переменные.
\end{itemize}

В каждый момент времени внутри монитора может выполняться не более одного процесса.

Условные переменные поддерживают операции:
\[
wait(c), \quad signal(c).
\]

Операция $wait(c)$ освобождает монитор и блокирует процесс до сигнала по условию $c$.

Пример монитора для ограниченного буфера:

\[
\begin{array}{l}
\textbf{monitor } BoundedBuffer\\
\quad buffer[N];\\
\quad count := 0;\\
\quad condition notFull, notEmpty;\\
\\
\quad \textbf{procedure } put(item):\\
\quad\quad \textbf{while } count = N \textbf{ do } wait(notFull);\\
\quad\quad insert(item);\\
\quad\quad count := count + 1;\\
\quad\quad signal(notEmpty);\\
\\
\quad \textbf{procedure } get():\\
\quad\quad \textbf{while } count = 0 \textbf{ do } wait(notEmpty);\\
\quad\quad item := remove();\\
\quad\quad count := count - 1;\\
\quad\quad signal(notFull);\\
\quad\quad \textbf{return } item.
\end{array}
\]

Использование цикла \texttt{while}, а не \texttt{if}, важно: после пробуждения условие может снова оказаться ложным.

\section{Межпроцессное взаимодействие через обмен сообщениями}

\subsection{Модель обмена сообщениями}

В модели \textbf{message passing} процессы не имеют общей памяти. Они взаимодействуют путём отправки и получения сообщений.

Базовые операции:
\[
send(destination, message),
\]
\[
receive(source, message).
\]

Общая схема:
\[
P_1 \xrightarrow{message} P_2.
\]

Сообщение может содержать:
\begin{itemize}
    \item данные;
    \item команду;
    \item идентификатор отправителя;
    \item временную метку;
    \item номер последовательности;
    \item служебную информацию.
\end{itemize}

\subsection{Преимущества обмена сообщениями}

\begin{itemize}
    \item естественно подходит для распределённых систем;
    \item отсутствие общей памяти снижает риск гонок;
    \item проще изолировать процессы друг от друга;
    \item коммуникация явно выражена в программе;
    \item хорошо сочетается с моделью акторов.
\end{itemize}

\subsection{Недостатки обмена сообщениями}

\begin{itemize}
    \item накладные расходы на копирование и сериализацию данных;
    \item сложность проектирования протоколов;
    \item возможность потери, дублирования или переупорядочивания сообщений в сетях;
    \item необходимость учитывать задержки;
    \item сложность обнаружения глобального состояния.
\end{itemize}

\subsection{Прямой и косвенный обмен сообщениями}

\subsubsection{Прямой обмен}

При прямом обмене отправитель явно указывает получателя:
\[
send(P_j, m),
\]
\[
receive(P_i, m).
\]

Свойства:
\begin{itemize}
    \item связь обычно устанавливается между парой процессов;
    \item процесс должен знать идентификатор партнёра;
    \item коммуникация менее гибкая при динамической структуре системы.
\end{itemize}

\subsubsection{Косвенный обмен}

При косвенном обмене сообщения отправляются не конкретному процессу, а в промежуточный объект --- \textbf{почтовый ящик}, очередь или канал.

\[
send(MB, m),
\]
\[
receive(MB, m).
\]

Здесь $MB$ --- mailbox.

Преимущества:
\begin{itemize}
    \item отправитель и получатель слабо связаны;
    \item несколько процессов могут читать из одного ящика;
    \item проще строить очереди задач;
    \item можно менять получателей без изменения отправителей.
\end{itemize}

\subsection{Синхронный и асинхронный обмен}

\subsubsection{Синхронная отправка}

При синхронной отправке операция:
\[
send(P, m)
\]
блокирует отправителя до тех пор, пока получатель не примет сообщение.

Такая схема называется \textbf{рандеву}:
\[
P_1 \text{ ждёт } P_2,
\quad
P_2 \text{ ждёт сообщение от } P_1.
\]

Преимущества:
\begin{itemize}
    \item сильная синхронизация;
    \item не требуется буферизация;
    \item отправитель знает, что сообщение принято.
\end{itemize}

Недостатки:
\begin{itemize}
    \item возможны блокировки;
    \item меньше параллелизма;
    \item хуже подходит для высоконагруженных систем.
\end{itemize}

\subsubsection{Асинхронная отправка}

При асинхронной отправке:
\[
send(P, m)
\]
возвращает управление сразу после помещения сообщения в буфер.

Преимущества:
\begin{itemize}
    \item выше степень параллелизма;
    \item отправитель не ждёт получателя;
    \item удобно для событийных систем.
\end{itemize}

Недостатки:
\begin{itemize}
    \item нужен буфер;
    \item возможна переполненность очереди;
    \item сложнее контролировать доставку.
\end{itemize}

\subsubsection{Блокирующее и неблокирующее получение}

Операция получения может быть:
\begin{itemize}
    \item \textbf{блокирующей}: процесс ждёт появления сообщения;
    \item \textbf{неблокирующей}: если сообщения нет, операция сразу возвращает специальный результат.
\end{itemize}

\[
receive(MB, m):
\begin{cases}
m, & \text{если сообщение есть},\\
\bot, & \text{если сообщения нет и операция неблокирующая}.
\end{cases}
\]

\subsection{Буферизация сообщений}

Канал связи может иметь:
\begin{enumerate}
    \item \textbf{нулевую ёмкость};
    \item \textbf{ограниченную ёмкость};
    \item \textbf{неограниченную ёмкость}.
\end{enumerate}

\subsubsection{Канал нулевой ёмкости}

Сообщение не хранится в буфере. Отправитель и получатель должны встретиться одновременно.

\[
capacity = 0.
\]

Это синхронное взаимодействие.

\subsubsection{Канал ограниченной ёмкости}

Буфер содержит не более $N$ сообщений:
\[
0 \leq count \leq N.
\]

Если буфер полон, отправитель блокируется или получает ошибку.

\subsubsection{Канал неограниченной ёмкости}

Теоретически буфер может расти неограниченно:
\[
capacity = \infty.
\]

На практике память всегда конечна, поэтому такая модель является абстракцией.

\subsection{Порядок доставки сообщений}

Возможные гарантии:
\begin{itemize}
    \item \textbf{FIFO-порядок}: сообщения от одного отправителя к одному получателю доставляются в порядке отправки;
    \item \textbf{каузальный порядок}: если отправка одного сообщения причинно предшествует другому, порядок доставки сохраняется;
    \item \textbf{полный порядок}: все процессы видят сообщения в одном и том же порядке.
\end{itemize}

\section{Организация почтовых ящиков}

\subsection{Понятие почтового ящика}

\textbf{Почтовый ящик} --- объект межпроцессного взаимодействия, представляющий логическую очередь сообщений.

Процессы могут выполнять операции:
\[
send(mailbox, message),
\]
\[
receive(mailbox, message).
\]

Графически:

\[
P_1, P_2, \ldots, P_k
\rightarrow
MB
\rightarrow
Q_1, Q_2, \ldots, Q_l.
\]

Здесь несколько отправителей помещают сообщения в ящик, а несколько получателей извлекают их.

\subsection{Свойства почтового ящика}

Почтовый ящик характеризуется:
\begin{itemize}
    \item именем или идентификатором;
    \item владельцем;
    \item правами доступа;
    \item ёмкостью;
    \item политикой очереди;
    \item способом блокировки процессов;
    \item временем жизни;
    \item надёжностью доставки.
\end{itemize}

\subsection{Виды почтовых ящиков}

\subsubsection{Системные почтовые ящики}

Создаются и управляются операционной системой. Обычно имеют системный идентификатор, права доступа и механизмы блокировки.

\subsubsection{Пользовательские почтовые ящики}

Реализуются на уровне приложения, например как структура данных:
\[
MB = (Queue, Mutex, NotEmpty, NotFull).
\]

\subsubsection{Именованные и неименованные ящики}

\textbf{Именованный} ящик доступен по имени:
\[
send("jobs", m).
\]

\textbf{Неименованный} ящик обычно передаётся процессу как дескриптор или ссылка.

\subsubsection{Одноадресные и многоадресные ящики}

В одноадресном ящике каждое сообщение получает один получатель.

В многоадресной модели сообщение может быть доставлено нескольким получателям:
\[
m \rightarrow \{P_1, P_2, \ldots, P_k\}.
\]

\subsection{Операции над почтовым ящиком}

Основные операции:
\begin{enumerate}
    \item создание:
    \[
    create\_mailbox(name, capacity);
    \]
    \item удаление:
    \[
    delete\_mailbox(name);
    \]
    \item отправка:
    \[
    send(name, message);
    \]
    \item получение:
    \[
    receive(name);
    \]
    \item проверка состояния:
    \[
    size(name),\quad isEmpty(name),\quad isFull(name).
    \]
\end{enumerate}

\subsection{Алгоритм работы ограниченного почтового ящика}

Пусть почтовый ящик имеет ёмкость $N$ и реализован очередью FIFO.

Состояние:
\[
Queue,\quad count,\quad N.
\]

Синхронизационные объекты:
\[
mutex = 1,\quad empty = N,\quad full = 0.
\]

Отправка сообщения:
\[
\begin{array}{l}
\textbf{procedure } send(MB, message):\\
\quad wait(empty);\\
\quad wait(mutex);\\
\quad enqueue(MB.Queue, message);\\
\quad MB.count := MB.count + 1;\\
\quad signal(mutex);\\
\quad signal(full).
\end{array}
\]

Получение сообщения:
\[
\begin{array}{l}
\textbf{procedure } receive(MB):\\
\quad wait(full);\\
\quad wait(mutex);\\
\quad message := dequeue(MB.Queue);\\
\quad MB.count := MB.count - 1;\\
\quad signal(mutex);\\
\quad signal(empty);\\
\quad \textbf{return } message.
\end{array}
\]

\subsubsection{Иллюстрация}

Пусть:
\[
N = 3,
\quad
Queue = [\,].
\]

После отправки сообщений $a,b,c$:
\[
Queue = [a,b,c],
\quad
count = 3.
\]

Следующая операция \texttt{send} блокируется, так как:
\[
empty = 0.
\]

После операции \texttt{receive}:
\[
Queue = [b,c],
\quad
count = 2,
\quad
empty = 1.
\]

Теперь один отправитель может быть разблокирован.

\subsection{Политики извлечения сообщений}

Почтовый ящик может использовать разные политики:
\begin{itemize}
    \item \textbf{FIFO}: первое отправленное сообщение извлекается первым;
    \item \textbf{LIFO}: последнее отправленное сообщение извлекается первым;
    \item \textbf{приоритетная очередь}: сообщения с большим приоритетом извлекаются раньше;
    \item \textbf{по шаблону}: получатель выбирает сообщение, удовлетворяющее условию;
    \item \textbf{по отправителю}: получатель ждёт сообщение от конкретного процесса.
\end{itemize}

\subsection{Почтовые ящики и модель акторов}

В модели акторов каждая активная сущность имеет собственный почтовый ящик. Актор:
\begin{itemize}
    \item получает сообщение из ящика;
    \item изменяет своё локальное состояние;
    \item отправляет сообщения другим акторам;
    \item может создавать новых акторов.
\end{itemize}

Акторы не имеют общей памяти. Это упрощает рассуждение о локальном состоянии, но переносит сложность в протоколы обмена сообщениями.

\section{Согласованность данных}

\subsection{Проблема согласованности}

В параллельных и распределённых системах одна и та же логическая переменная или объект может иметь несколько копий:
\begin{itemize}
    \item в кэшах процессоров;
    \item в оперативной памяти разных узлов;
    \item в репликах базы данных;
    \item в локальных буферах процессов;
    \item в распределённых хранилищах.
\end{itemize}

\textbf{Согласованность данных} --- это набор гарантий о том, какие значения могут видеть операции чтения при наличии параллельных операций записи.

Проблема особенно важна для:
\begin{itemize}
    \item многопроцессорных архитектур;
    \item распределённых баз данных;
    \item файловых систем;
    \item кэшей;
    \item реплицированных хранилищ;
    \item распределённых объектов.
\end{itemize}

\subsection{Операции чтения и записи}

Обычно рассматривают объекты, над которыми выполняются операции:
\[
read(x) \rightarrow v,
\]
\[
write(x, v).
\]

Операция записи помещает значение $v$ в объект $x$.

Операция чтения возвращает некоторое значение $v$ объекта $x$.

В последовательной программе естественно ожидать, что чтение возвращает последнее записанное значение. В параллельной системе понятие ``последнее'' становится неоднозначным.

\subsection{Истории выполнения}

\textbf{История} --- последовательность событий вызова и завершения операций над разделяемыми объектами.

Например:
\[
write(x,1),\quad read(x)\rightarrow 1,\quad write(x,2),\quad read(x)\rightarrow 2.
\]

В параллельном случае операции могут перекрываться во времени:

\[
\begin{array}{c|c|c}
\text{Время} & P_1 & P_2 \\
\hline
1 & write(x,1)\ \text{начало} & \\
2 & & read(x)\ \text{начало} \\
3 & write(x,1)\ \text{конец} & \\
4 & & read(x)\rightarrow ?
\end{array}
\]

Вопрос: какое значение допустимо вернуть?

Ответ зависит от модели согласованности.

\section{Модели согласованности данных}

\subsection{Строгая согласованность}

\textbf{Строгая согласованность} требует, чтобы любая операция чтения возвращала результат самой последней операции записи в абсолютном глобальном времени.

Формально:
\[
read(x) \rightarrow v
\]
должна вернуть значение последней операции:
\[
write(x,v),
\]
которая завершилась до момента чтения в глобальном времени.

Эта модель интуитивно наиболее проста, но практически недостижима в распределённых системах, поскольку требует:
\begin{itemize}
    \item точных глобальных часов;
    \item мгновенного распространения записей;
    \item нулевых задержек коммуникации.
\end{itemize}

\subsection{Последовательная согласованность}

\textbf{Последовательная согласованность} была сформулирована Лесли Лэмпортом.

\textbf{Определение.}
Система является последовательно согласованной, если результат любого выполнения такой же, как если бы операции всех процессов были выполнены в некотором последовательном порядке, причём операции каждого отдельного процесса в этом порядке идут в том же порядке, в котором они были заданы этим процессом.

Иными словами, должна существовать глобальная последовательность операций $S$, такая что:
\begin{enumerate}
    \item $S$ содержит все операции всех процессов;
    \item порядок операций каждого процесса в $S$ сохраняется;
    \item каждое чтение возвращает значение последней предшествующей записи в $S$.
\end{enumerate}

Важно: последовательная согласованность не требует сохранения реального времени между операциями разных процессов.

\subsubsection{Пример последовательной согласованности}

Пусть изначально:
\[
x = 0,\quad y = 0.
\]

Два процесса:
\[
P_1: write(x,1); \quad r_1 := read(y),
\]
\[
P_2: write(y,1); \quad r_2 := read(x).
\]

Результат:
\[
r_1 = 0,\quad r_2 = 0
\]
может быть допустим в некоторых слабых моделях, но не является последовательно согласованным.

\textbf{Почему?}

Чтобы $r_1 = 0$, чтение $read(y)$ в $P_1$ должно стоять до $write(y,1)$ в глобальном порядке:
\[
read(y) < write(y,1).
\]

Но порядок процесса $P_1$ требует:
\[
write(x,1) < read(y).
\]

Чтобы $r_2 = 0$, чтение $read(x)$ в $P_2$ должно стоять до $write(x,1)$:
\[
read(x) < write(x,1).
\]

Порядок процесса $P_2$ требует:
\[
write(y,1) < read(x).
\]

Итого:
\[
write(x,1) < read(y) < write(y,1) < read(x) < write(x,1).
\]

Получен цикл:
\[
write(x,1) < write(x,1),
\]
что невозможно. Следовательно, такой результат не является последовательно согласованным.

\subsubsection{Теорема: строгая согласованность влечёт последовательную}

\textbf{Теорема.}
Если система строго согласована, то она последовательно согласована.

\textbf{Доказательство.}

Пусть система строго согласована. Тогда каждое чтение возвращает значение последней записи в абсолютном глобальном времени.

Рассмотрим произвольное выполнение системы. Упорядочим все операции по моментам их выполнения в глобальном времени. Если операции не перекрываются, порядок задаётся непосредственно. Если операции перекрываются, выберем любой порядок, согласованный с моментами завершения и начала так, чтобы не нарушать порядок операций внутри каждого процесса.

Получим некоторую глобальную последовательность $S$ всех операций.

Поскольку в строгой модели каждое чтение возвращает последнюю запись по глобальному времени, в последовательности $S$ это же чтение возвращает последнюю предшествующую запись. Кроме того, операции каждого процесса в $S$ идут в исходном программном порядке, так как выполнение одного процесса само упорядочено во времени.

Следовательно, существует глобальная последовательность операций, сохраняющая порядок каждого процесса и объясняющая результаты всех чтений. Это ровно условие последовательной согласованности.

Теорема доказана.
\[
\blacksquare
\]

\subsection{Линеаризуемость}

\textbf{Линеаризуемость} --- более сильная модель согласованности для операций над объектами.

\textbf{Определение.}
История операций является линеаризуемой, если каждой операции можно сопоставить мгновенную точку выполнения, называемую \textbf{точкой линеаризации}, между её вызовом и ответом так, что получившаяся последовательная история:
\begin{enumerate}
    \item корректна для соответствующего последовательного объекта;
    \item сохраняет реальный порядок неперекрывающихся операций.
\end{enumerate}

Если операция $A$ завершилась до начала операции $B$, то в линеаризованном порядке:
\[
A < B.
\]

Линеаризуемость сильнее последовательной согласованности, потому что учитывает реальное время.

\subsubsection{Пример}

Пусть изначально:
\[
x = 0.
\]

История:
\[
P_1: write(x,1) \text{ завершилась},
\]
затем после её завершения:
\[
P_2: read(x) \rightarrow 0.
\]

Такая история не линеаризуема, потому что чтение началось после завершения записи и должно видеть $1$.

Она также не будет строгой. Однако последовательная согласованность в некоторых формулировках могла бы не учитывать реальное время между разными процессами, если сохраняется только программный порядок.

\subsubsection{Теорема: линеаризуемость влечёт последовательную согласованность}

\textbf{Теорема.}
Всякая линеаризуемая история является последовательно согласованной.

\textbf{Доказательство.}

Пусть история $H$ линеаризуема. По определению, существует последовательная история $S$, полученная выбором точки линеаризации для каждой операции внутри интервала между её вызовом и ответом.

Эта последовательная история:
\begin{enumerate}
    \item содержит все операции из $H$;
    \item корректна относительно последовательной спецификации объектов;
    \item сохраняет реальный порядок неперекрывающихся операций.
\end{enumerate}

Операции одного процесса не перекрываются с точки зрения этого процесса: процесс вызывает следующую операцию после завершения предыдущей. Поэтому реальный порядок операций каждого процесса сохраняется в $S$.

Следовательно, существует последовательный порядок всех операций, сохраняющий программный порядок каждого процесса и объясняющий возвращаемые значения. Это и есть условие последовательной согласованности.

Теорема доказана.
\[
\blacksquare
\]

\subsection{Причинная согласованность}

\textbf{Причинная согласованность} учитывает отношение причинности между операциями.

Операция $a$ причинно предшествует операции $b$, если:
\begin{enumerate}
    \item $a$ и $b$ выполнены одним процессом, и $a$ была раньше $b$;
    \item $a$ была записью, значение которой прочитала операция $b$;
    \item существует цепочка причинных зависимостей:
    \[
    a \rightarrow c \rightarrow b.
    \]
\end{enumerate}

Обозначают:
\[
a \rightarrow b.
\]

\textbf{Определение.}
Система причинно согласована, если все процессы видят причинно связанные записи в одном и том же порядке. Конкурентные записи могут наблюдаться разными процессами в разном порядке.

\subsubsection{Пример}

Пусть:
\[
P_1: write(x,1).
\]
Затем процесс $P_2$ читает:
\[
read(x) \rightarrow 1
\]
и после этого пишет:
\[
write(y,1).
\]

Тогда:
\[
write(x,1) \rightarrow write(y,1).
\]

Все процессы, которые видят обе записи, должны видеть сначала:
\[
write(x,1),
\]
а затем:
\[
write(y,1).
\]

Но если две записи независимы:
\[
P_1: write(x,1),
\]
\[
P_2: write(y,1),
\]
то разные процессы могут увидеть их в разном порядке.

\subsection{FIFO-согласованность}

\textbf{FIFO-согласованность}, или PRAM-согласованность, требует, чтобы записи каждого отдельного процесса наблюдались всеми процессами в порядке их выполнения этим процессом.

Если:
\[
P_i: write(x,1); write(y,1),
\]
то все процессы должны увидеть запись $x:=1$ раньше записи $y:=1$.

Но записи разных процессов могут наблюдаться в разных порядках.

Отношение между моделями:
\[
\text{последовательная согласованность}
\Rightarrow
\text{причинная согласованность}
\Rightarrow
\text{FIFO-согласованность}.
\]

\subsection{Слабая согласованность}

\textbf{Слабая согласованность} разделяет обычные операции доступа к данным и операции синхронизации.

Идея: не требуется немедленно согласовывать все записи. Достаточно гарантировать согласованность в точках синхронизации.

Типичные правила:
\begin{enumerate}
    \item доступ к синхронизационной переменной является последовательно согласованным;
    \item перед выполнением операции синхронизации должны быть завершены предыдущие записи;
    \item обычные чтения и записи могут переупорядочиваться между синхронизациями.
\end{enumerate}

\subsection{Release consistency}

\textbf{Release consistency} уточняет слабую согласованность, разделяя синхронизационные операции на:
\begin{itemize}
    \item \textbf{acquire} --- захват ресурса;
    \item \textbf{release} --- освобождение ресурса.
\end{itemize}

Операция:
\[
acquire(L)
\]
должна обеспечить видимость данных, защищённых синхронизационным объектом $L$.

Операция:
\[
release(L)
\]
должна распространить изменения, сделанные в критической секции.

Типичная схема:
\[
acquire(L);
\]
\[
\text{работа с разделяемыми данными;}
\]
\[
release(L).
\]

Смысл:
\begin{itemize}
    \item перед \texttt{release} все предыдущие записи должны стать доступными другим процессам, которые затем выполнят \texttt{acquire};
    \item после \texttt{acquire} процесс должен видеть изменения, опубликованные соответствующим \texttt{release}.
\end{itemize}

\subsection{Eventual consistency}

\textbf{Согласованность в конечном счёте} используется в распределённых хранилищах и реплицированных системах.

\textbf{Определение.}
Если новые обновления объекта прекращаются, то через конечное время все реплики объекта сойдутся к одному и тому же значению.

Иначе говоря:
\[
\text{no new writes} \Rightarrow \exists T: \forall t > T,\ replicas(t) \text{ equal}.
\]

Преимущества:
\begin{itemize}
    \item высокая доступность;
    \item устойчивость к сетевым разделениям;
    \item хорошая масштабируемость.
\end{itemize}

Недостатки:
\begin{itemize}
    \item чтение может вернуть устаревшие данные;
    \item возможны конфликты обновлений;
    \item приложению часто приходится разрешать конфликты.
\end{itemize}

\subsection{Read-your-writes consistency}

\textbf{Read-your-writes consistency} гарантирует, что процесс всегда видит собственные записи.

Если процесс выполнил:
\[
write(x,1),
\]
то последующее чтение этим же процессом:
\[
read(x)
\]
не должно вернуть старое значение, например $0$.

Эта гарантия важна для пользовательских приложений. Например, если пользователь изменил профиль, он ожидает сразу видеть новое значение.

\subsection{Monotonic reads}

\textbf{Monotonic reads} означает, что если процесс однажды увидел некоторое значение объекта, то позже он не должен увидеть более старое значение.

Если:
\[
read(x) \rightarrow v_2
\]
после того, как уже было увидено:
\[
read(x) \rightarrow v_1,
\]
и $v_2$ старее $v_1$, то это нарушение monotonic reads.

\subsection{Monotonic writes}

\textbf{Monotonic writes} требует, чтобы записи одного процесса применялись в порядке их выполнения.

Если процесс выполнил:
\[
write(x,1);
\]
затем:
\[
write(x,2),
\]
то система не должна применить их как:
\[
write(x,2) \text{ раньше } write(x,1).
\]

\subsection{Writes-follow-reads}

\textbf{Writes-follow-reads} требует, чтобы запись, выполненная после чтения некоторого значения, была упорядочена после записи, породившей это прочитанное значение.

Например:
\[
read(x) \rightarrow 1;
\]
затем:
\[
write(y,1).
\]
Тогда запись $y:=1$ должна причинно следовать за записью, которая установила $x:=1$.

\section{Связь согласованности с синхронизацией}

\subsection{Почему синхронизация важна для согласованности}

В современных многопроцессорных системах процессор и компилятор могут переупорядочивать операции для оптимизации.

Например, код:
\[
x := 1;
\]
\[
ready := true;
\]
может быть виден другому потоку так, будто сначала изменился \texttt{ready}, а потом \texttt{x}, если отсутствуют барьеры памяти.

Другой поток:
\[
\textbf{while } ready = false \textbf{ do wait;}
\]
\[
print(x);
\]

может теоретически увидеть:
\[
ready = true,
\quad
x = 0.
\]

Для предотвращения таких эффектов применяют:
\begin{itemize}
    \item мьютексы;
    \item атомарные переменные;
    \item барьеры памяти;
    \item acquire/release-семантику;
    \item volatile-переменные в некоторых языках, но их семантика зависит от языка.
\end{itemize}

\subsection{Атомарные операции}

\textbf{Атомарная операция} выполняется как неделимое действие относительно других потоков.

Примеры:
\begin{itemize}
    \item test-and-set;
    \item compare-and-swap;
    \item fetch-and-add;
    \item load-linked/store-conditional.
\end{itemize}

\subsection{Compare-and-swap}

Операция \texttt{CAS} имеет вид:
\[
CAS(addr, expected, new).
\]

Семантика:
\[
\begin{array}{l}
\textbf{if } *addr = expected \textbf{ then}\\
\quad *addr := new;\\
\quad \textbf{return } true;\\
\textbf{else}\\
\quad \textbf{return } false.
\end{array}
\]

CAS используется для построения неблокирующих структур данных.

Минимальный пример: атомарный инкремент.

\[
\begin{array}{l}
\textbf{procedure } atomic\_increment(x):\\
\quad \textbf{repeat}\\
\quad\quad old := x;\\
\quad\quad new := old + 1;\\
\quad \textbf{until } CAS(x, old, new) = true.
\end{array}
\]

Идея: если между чтением \texttt{old} и попыткой записи другой поток изменил $x$, операция CAS завершится неудачей, и попытка повторится.

\section{Взаимные блокировки}

\subsection{Определение deadlock}

\textbf{Взаимная блокировка} возникает, когда множество процессов находится в состоянии ожидания событий, которые могут быть вызваны только процессами из того же множества.

Пример:
\[
P_1 \text{ удерживает } R_1 \text{ и ждёт } R_2,
\]
\[
P_2 \text{ удерживает } R_2 \text{ и ждёт } R_1.
\]

Граф ожидания:
\[
P_1 \rightarrow R_2 \rightarrow P_2 \rightarrow R_1 \rightarrow P_1.
\]

\subsection{Условия Коффмана}

Для возникновения deadlock необходимы четыре условия:

\begin{enumerate}
    \item \textbf{Взаимное исключение}: ресурс не может одновременно использоваться несколькими процессами.
    \item \textbf{Удержание и ожидание}: процесс удерживает один ресурс и ждёт другой.
    \item \textbf{Отсутствие вытеснения}: ресурс нельзя принудительно отобрать.
    \item \textbf{Циклическое ожидание}: существует цикл процессов, каждый из которых ждёт ресурс, удерживаемый следующим.
\end{enumerate}

\subsubsection{Теорема о необходимости условий Коффмана}

\textbf{Теорема.}
Если в системе возникла взаимная блокировка, то выполняются все четыре условия Коффмана.

\textbf{Доказательство.}

Пусть в системе возникла взаимная блокировка.

\textbf{1. Взаимное исключение.}
Если бы все ресурсы могли использоваться совместно неограниченным числом процессов, то процесс не мог бы блокироваться из-за занятости ресурса. Следовательно, для deadlock необходим хотя бы один ресурс с режимом взаимного исключения.

\textbf{2. Удержание и ожидание.}
В состоянии deadlock процесс не просто ждёт ресурс, но при этом препятствует продвижению других процессов, удерживая ресурсы. Если бы процессы не удерживали ресурсы во время ожидания, ресурсы могли бы быть освобождены и переданы ожидающим.

\textbf{3. Отсутствие вытеснения.}
Если бы ресурс можно было принудительно отобрать у процесса, то система могла бы разорвать ожидание, передав ресурс другому процессу. Deadlock сохраняется именно потому, что ресурсы нельзя отобрать без согласия владельца.

\textbf{4. Циклическое ожидание.}
В deadlock каждый процесс ждёт событие, которое может быть вызвано другим заблокированным процессом. Если построить граф ожидания, где ребро $P_i \rightarrow P_j$ означает, что $P_i$ ждёт ресурс, удерживаемый $P_j$, то из конечности множества заблокированных процессов и отсутствия возможности продвижения следует наличие цикла ожидания.

Значит, все четыре условия необходимы.

Теорема доказана.
\[
\blacksquare
\]

\subsection{Методы борьбы с deadlock}

\begin{itemize}
    \item \textbf{предотвращение} --- запрещение хотя бы одного из условий Коффмана;
    \item \textbf{избежание} --- динамический анализ безопасных состояний;
    \item \textbf{обнаружение и восстановление} --- разрешить deadlock, затем обнаружить и устранить;
    \item \textbf{игнорирование} --- подход, применяемый в ряде ОС, если deadlock редок и цена предотвращения высока.
\end{itemize}

\subsection{Упорядочение ресурсов}

Простой способ предотвращения deadlock --- назначить ресурсам глобальный порядок:
\[
R_1 < R_2 < \cdots < R_n.
\]

Процесс может запрашивать ресурсы только в возрастающем порядке.

\subsubsection{Теорема}

\textbf{Теорема.}
Если все процессы запрашивают ресурсы только в порядке возрастания глобальной нумерации, то циклическое ожидание невозможно.

\textbf{Доказательство.}

Предположим противное: существует цикл ожидания:
\[
P_1 \rightarrow P_2 \rightarrow \cdots \rightarrow P_k \rightarrow P_1.
\]

Пусть процесс $P_i$ удерживает ресурс $R_i$ и ждёт ресурс $R_{i+1}$, удерживаемый $P_{i+1}$.

Так как процессы запрашивают ресурсы только в возрастающем порядке, из того, что $P_i$ удерживает $R_i$ и ждёт $R_{i+1}$, следует:
\[
R_i < R_{i+1}.
\]

Следовательно:
\[
R_1 < R_2 < \cdots < R_k < R_1.
\]

Получаем:
\[
R_1 < R_1,
\]
что невозможно.

Противоречие. Значит, циклическое ожидание невозможно.

Теорема доказана.
\[
\blacksquare
\]

\section{Сравнение общей памяти и обмена сообщениями}

\[
\begin{array}{|p{4cm}|p{5cm}|p{5cm}|}
\hline
\textbf{Критерий} & \textbf{Общая память} & \textbf{Обмен сообщениями} \\
\hline
Способ взаимодействия & Чтение и запись разделяемых данных & Отправка и получение сообщений \\
\hline
Синхронизация & Требуется явно: мьютексы, семафоры, мониторы & Частично встроена в операции send/receive \\
\hline
Производительность & Высокая на одной машине & Зависит от стоимости передачи сообщений \\
\hline
Распределённость & Плохо подходит напрямую & Естественно подходит \\
\hline
Опасность гонок & Высокая & Ниже, если нет общей памяти \\
\hline
Отладка & Часто сложная & Протоколы сложны, но обмен явно виден \\
\hline
Масштабирование & Ограничено кэш-когерентностью и памятью & Хорошо масштабируется в распределённых системах \\
\hline
Типичные примеры & Потоки, shared memory, memory-mapped files & MPI, сокеты, очереди сообщений, акторы \\
\hline
\end{array}
\]

\section{Минимальные примеры}

\subsection{Пример взаимодействия через общую память}

Пусть два процесса должны передать число через общую переменную.

Разделяемые переменные:
\[
data,\quad ready = false.
\]

Производитель:
\[
\begin{array}{l}
data := 42;\\
ready := true.
\end{array}
\]

Потребитель:
\[
\begin{array}{l}
\textbf{while } ready = false \textbf{ do wait};\\
print(data).
\end{array}
\]

Этот пример некорректен без синхронизации в современных системах, потому что возможны:
\begin{itemize}
    \item переупорядочивание операций;
    \item кэширование;
    \item отсутствие гарантии видимости.
\end{itemize}

Корректный вариант с мьютексом и условной переменной:
\[
\begin{array}{l}
\textbf{producer:}\\
\quad lock(m);\\
\quad data := 42;\\
\quad ready := true;\\
\quad signal(c);\\
\quad unlock(m).
\end{array}
\]

\[
\begin{array}{l}
\textbf{consumer:}\\
\quad lock(m);\\
\quad \textbf{while } ready = false \textbf{ do } wait(c,m);\\
\quad print(data);\\
\quad unlock(m).
\end{array}
\]

\subsection{Пример взаимодействия через сообщения}

Процесс $P_1$:
\[
send(P_2, 42).
\]

Процесс $P_2$:
\[
x := receive(P_1).
\]

Если \texttt{send} и \texttt{receive} синхронные, то $P_1$ будет ждать, пока $P_2$ не выполнит получение.

Если обмен асинхронный, сообщение сначала помещается в буфер:
\[
P_1 \rightarrow Buffer \rightarrow P_2.
\]

\subsection{Пример почтового ящика как очереди задач}

Пусть есть почтовый ящик:
\[
jobs.
\]

Master:
\[
\begin{array}{l}
\textbf{for } i := 1 \textbf{ to } n \textbf{ do}\\
\quad send(jobs, task_i).
\end{array}
\]

Worker:
\[
\begin{array}{l}
\textbf{while true do}\\
\quad task := receive(jobs);\\
\quad result := process(task);\\
\quad send(results, result).
\end{array}
\]

Эта схема реализует динамическое распределение нагрузки: свободный worker сам забирает следующую задачу.

\section{Краткая сводка}

Параллельные процессы позволяют выполнять несколько вычислений с перекрытием во времени. Они могут порождаться статически, динамически, иерархически, по схемам fork--join, master--worker или конвейера.

Управление параллельными процессами бывает централизованным и децентрализованным, синхронным и асинхронным.

Межпроцессное взаимодействие организуется двумя фундаментальными способами:
\begin{itemize}
    \item через общую память;
    \item через обмен сообщениями.
\end{itemize}

Общая память эффективна, но требует строгой синхронизации. Обмен сообщениями лучше подходит для распределённых систем и явно выражает коммуникацию, но требует проектирования протоколов и буферизации.

Почтовые ящики являются формой косвенного обмена сообщениями. Они позволяют отделить отправителя от получателя и естественно реализуют очереди задач, акторные системы и асинхронное взаимодействие.

Модели согласованности определяют, какие значения могут возвращать операции чтения в параллельной или распределённой системе. Наиболее важные модели:
\begin{itemize}
    \item строгая согласованность;
    \item линеаризуемость;
    \item последовательная согласованность;
    \item причинная согласованность;
    \item FIFO-согласованность;
    \item слабая согласованность;
    \item release consistency;
    \item eventual consistency.
\end{itemize}

Чем сильнее модель согласованности, тем проще рассуждать о корректности программы, но тем выше стоимость реализации и ниже потенциальная масштабируемость.

\end{document}
